% !PRINTED 552
\weberpage{552}
%<<CONT
sei, worin $a$ eine konstante (unabhängig von $t$), und $\theta$ eine Funktion von $t$, die in endlichen Grenzen eingeschlossen bleibt.

Nehmen wir zunächst an, die $\mu_n$ seien alle voneinander verschieden, so ist $Z(\mu_n) = n$, und aus (2) folgt:
\weq{3}{\frac{n}{\mu_n} = a + \frac{\theta}{\sqrt{\mu_n}};}
daraus folgt, daß $n:\mu_n$ endlich bleibt, und folglich mit veränderter Bedeutung von $\theta$:
\[
\frac{n}{\mu_n} = a + \frac{\theta}{\sqrt{n}},
\]
oder, nach dem binomischen Satz, für irgend ein positives $s$:
\weq{4}{\frac{n^s}{\mu_n^s} = a^s\left(1 + \frac{\theta}{\sqrt{n}}\right),}
worin alle die mit $\theta$ bezeichneten Größen endliche Werte haben.

Diese Formeln gelten aber auch noch, wenn unter den $\mu_n$ gleiche vorkommen. Denn seien etwa (wie in Bd.\,II, § 196, 3.)
\[
\mu_{m+1}, \mu_{m+2}, \ldots, \mu_{m+l} = \mu_n
\]
einander gleich und
\[
m < n \leqq m+l,
\]
so ist
\[
Z(\mu_n-0) = m, \qquad Z(\mu_n) = m+l,
\]
\[
\frac{Z(\mu_n-0)}{\mu_n} < \frac{n}{\mu_n} \leqq \frac{Z(\mu_n)}{\mu_n},
\]
\[
a + \frac{\theta'}{\sqrt{\mu_n}} < \frac{n}{\mu_n} < a + \frac{\theta}{\sqrt{\mu_n}},
\]
woraus (3) folgt und (4) wie oben abgeleitet werden kann.

Der Satz, den wir beweisen wollen, lautet so:

\gesperrt{Es sei $\varepsilon_1, \varepsilon_2, \varepsilon_3, \ldots$ eine unendliche Reihe positiver oder negativer, aber endlicher Größen und}
\weq{5}{\varepsilon_1+\varepsilon_2+\cdots+\varepsilon_n = \gamma_n\sqrt{n},}
\gesperrt{so beschaffen, daß $\gamma_n$ dem absoluten Werte nach unter einer endlichen Konstanten bleibt.}

\gesperrt{Dann ist}
\weq{6}{\sigma = \frac{\varepsilon_1}{\mu_1^s} + \frac{\varepsilon_2}{\mu_2^s} + \frac{\varepsilon_3}{\mu_3^s} + \cdots}
\gesperrt{konvergent und eine stetige Funktion von $s$, so lange}
\[
s > \tfrac{1}{2}
\]
\gesperrt{ist.}
% [sic] Der gedruckte binomische Schritt(4) setzt a>0 voraus; die formulierte Voraussetzung schliesst a=0 nicht aus. Kein Eingriff in die Formel.
